\section{Immutable Parent Pointers}

\subsection{Overview}
A cornerstone property of the Recursive Spawning framework is the immutability of parent-child relationships once established. Every spawned entity carries a reference to its parent, encoded as a \texttt{ParentPointer} that points from child to parent. This pointer, once assigned at the moment of spawn authorization, is permanently sealed and can never be retroactively modified, relocated, or nullified. The immutability guarantee is not a policy preference but a structural enforcement at the Fact Layer, making the spawn chain an immutable chronicle rather than a mutable ledger.

The immutability of parent pointers gives the spawn chain its append-only character. Because ancestry can never be rewritten, the chain serves as a reliable, tamper-evident record of causal provenance. This property is critical for downstream operations that depend on ancestry traversal — auditing, authorization, inheritance resolution, and liability attribution all require the guarantee that the chain reflects what actually happened, not what someone later preferred to have happened.

\subsection{ParentPointer: Definition and Assignment}

\begin{definition}[ParentPointer]
Given a parent entity $p$ and a child entity $c$, the \textbf{ParentPointer} of $c$, denoted $\texttt{ParentPointer}(c)$, is a reference to $p$. It is assigned exactly once at the moment $c$ is spawned and is then permanently sealed. $\texttt{ParentPointer}(c) = p$ is the sole admissible value for the lifetime of $c$.
\end{definition}

The assignment of a \texttt{ParentPointer} occurs during the spawn event itself, after the Purpose Layer has granted authorization (see \S4.4). The Fact Layer records the pointer as part of the entity's foundational state. There is no operation — administrative, judicial, or computational — that can alter this pointer after the fact. The immutability guarantee is enforced by the architecture of the Fact Layer, not by access-control policies that could be overridden.

For the root human in any spawn chain, the base case applies: $\texttt{ParentPointer}(root) = \bot$ (null), indicating that this entity has no parent. The null parent is not a placeholder but a deliberate and permanent property, indicating that the root human initiated their existence outside the spawn system and serves as the terminus of the chain in the upward direction.

\subsection{The Spawn Chain: Recursive Construction}

The spawn chain for any entity is constructed recursively by following parent pointers upward until the root is reached. This chain encodes the full ancestry of an entity in temporal order from root to self.

\begin{definition}[Chain]
For any entity $a$, the \textbf{Chain} of $a$, denoted $\texttt{Chain}(a)$, is defined recursively as:

\[
\texttt{Chain}(a) =
\begin{cases}
\{(a, \bot)\} & \text{if } \texttt{ParentPointer}(a) = \bot \text{ (root case)} \\[6pt]
\{(a, \texttt{ParentPointer}(a))\} \cup \texttt{Chain}(\texttt{ParentPointer}(a)) & \text{otherwise}
\end{cases}
\]

Each element $(child, parent)$ in the chain is an ordered pair representing a direct parent-child link. The chain is a set, preserving uniqueness, and is inherently ordered by the recursive construction from root to leaf.
\end{definition}

The base case terminates at the root human, where the parent pointer is null. Every non-root entity's chain contains the pair $(a, p)$ where $p = \texttt{ParentPointer}(a)$, followed by all pairs in $\texttt{Chain}(p)$. The recursive construction guarantees that the chain is complete, unambiguous, and free of cycles — by construction, following parent pointers always moves toward the root, never sideways or downward, because parent pointers are assigned at spawn and never modified.

\begin{algorithm}[Chain Traversal]
\textbf{Input:} An entity $a$ \\[4pt]
\textbf{Output:} The ordered set $\texttt{Chain}(a)$

\begin{enumerate}
    \item Initialize $current \gets a$, $chain \gets \emptyset$
    \item \textbf{repeat}:
    \begin{itemize}
        \item $parent \gets \texttt{ParentPointer}(current)$
        \item $chain \gets chain \cup \{(current, parent)\}$
        \item \textbf{if} $parent = \bot$ \textbf{then} \textbf{break}
        \item $current \gets parent$
    \end{itemize}
    \item \textbf{return} $chain$
\end{enumerate}

The algorithm terminates because each iteration assigns $current \gets parent$, strictly ascending the chain, and the chain is finite (terminates at the root where parent is $\bot$).
\end{algorithm}

The chain traversal algorithm has $O(d)$ time complexity where $d$ is the depth of the entity in the spawn tree. The depth of any entity is bounded by the total number of spawned entities in the system, making worst-case traversal linear in the number of entities in the chain.

\subsection{Immutability Enforcement at the Fact Layer}

The Fact Layer is the architectural locus of immutability enforcement. Unlike permission-based access control, which operates at the logic layer and can be reconfigured, the Fact Layer encodes immutability constraints that are enforced by the system's substrate. This means immutability is not a policy that can be changed by an administrator — it is a property of how the Fact Layer stores and retrieves parent pointers.

The enforcement mechanism operates as follows. When an entity is spawned, the Fact Layer records the entity's initial state, including its \texttt{ParentPointer}. This record is written to the substrate storage with an immutability flag that prohibits update operations on that field. The storage layer rejects any attempt to modify a \texttt{ParentPointer} field after write — including attempts by system administrators, Purpose Layer authorizers, or any other privileged component.

The rationale for substrate-level enforcement is to prevent the class of failures that arise from rely-once trusted computations. If immutability were enforced only at the application layer, a single misconfiguration, code defect, or compromised credential could allow parent pointer manipulation, invalidating the chain's integrity and undermining every downstream computation that depends on ancestry. By making immutability a property of the storage substrate, the system eliminates this entire failure class.

The Fact Layer also rejects \texttt{ParentPointer} reassignment during entity migration, replication, or state synchronization. No operation — local or distributed — can alter the pointer. This ensures that the chain is consistent across all nodes and replicas in the system. If a spawn chain is replicated for fault tolerance or geographic distribution, every replica carries bit-for-bit identical parent pointer values.

\subsection{Spawn Authorization at the Purpose Layer}

The Purpose Layer governs the conditions under which a new entity may be spawned. Authorization at this layer determines whether a parent may produce a child, whether the spawn event is consistent with the parent's purpose and the system's design intent, and whether the resulting child entity should be instantiated with a given set of initial properties.

When the Purpose Layer grants spawn authorization, it issues a directive to the Fact Layer to instantiate the new entity with the assigned \texttt{ParentPointer}. The Fact Layer does not evaluate authorization — it receives an authorization token and executes the instantiation. This separation of concerns ensures that authorization logic (Purpose Layer) is cleanly decoupled from factual record-keeping (Fact Layer).

The Purpose Layer's authorization decision may consider factors including but not limited to: the parent's remaining spawn quota, the alignment between the proposed child's purpose and the parent's registered intent, compliance with jurisdictional spawn regulations, and verification that the parent entity remains in an active, authorized state. These factors are evaluated against the purpose registry and policy constraints at the time of the request. The Purpose Layer does not, however, possess the ability to modify a \texttt{ParentPointer} after the Fact Layer has recorded it — authorization governs the creation event, not subsequent modification.

The spawn authorization process can be summarized as:
\begin{enumerate}
    \item Parent $p$ submits a spawn request to the Purpose Layer.
    \item The Purpose Layer evaluates the request against purpose registry and spawn policy.
    \item Upon authorization, the Purpose Layer issues a spawn directive to the Fact Layer.
    \item The Fact Layer instantiates child $c$ with $\texttt{ParentPointer}(c) = p$.
    \item The immutability seal is applied; no further modification is possible.
\end{enumerate}

This sequence ensures that every child entity's parentage is authorized at the moment of creation, and thereafter becomes permanently immutable.

\subsection{Properties of the Immutable Chain}

The immutability of parent pointers yields several important properties that propagate throughout the framework.

\textbf{Non-repudiation.} Because parent pointers cannot be altered after assignment, an entity's recorded ancestry is a definitive statement of its origin. An entity cannot later deny its spawn lineage, because the chain cannot be modified to reflect a different history. This property is critical for liability attribution, inheritance resolution, and audit trails.

\textbf{Cycle-freedom.} The chain is inherently acyclic by construction. Since $\texttt{ParentPointer}(c) = p$ is assigned at spawn and never changes, and each iteration of chain traversal moves from child to parent (strictly upward), there is no possible path that loops back to any prior entity. The spawn tree is a rooted arborescence with the root human as its apex.

\textbf{Auditability.} Any entity's complete ancestry can be reconstructed by a simple traversal of parent pointers. The traversal is deterministic, idempotent, and requires no external state — only the immutable records maintained by the Fact Layer. This makes ancestry audits inexpensive and tamper-evident.

\textbf{Temporal ordering.} The chain implicitly encodes temporal order because each parent must exist before its child is spawned. Following the chain from any entity to the root reconstructs the temporal sequence of the spawn events in order from newest to oldest. This ordering is fixed and cannot be rearranged.

\begin{figure}[htbp]
\centering
\begin{tikzpicture}[>=stealth, font=\small, every node/.style={minimum height=0.6cm, outer sep=2pt}]

% Nodes - Human root at top, then spawned entities below
\node[draw, circle, fill=blue!20, minimum width=1.2cm] (H) at (0, 0) {$H_0$};
\node[draw, circle, fill=green!20, minimum width=1.2cm] (A) at (-2.5, -2) {$A_1$};
\node[draw, circle, fill=green!20, minimum width=1.2cm] (B) at (2.5, -2) {$B_1$};
\node[draw, circle, fill=yellow!30, minimum width=1.2cm] (C) at (-3.5, -4) {$C_2$};
\node[draw, circle, fill=yellow!30, minimum width=1.2cm] (D) at (-1.5, -4) {$D_2$};
\node[draw, circle, fill=yellow!30, minimum width=1.2cm] (E) at (1.5, -4) {$E_2$};

% Arrows showing parent pointers (pointing from child to parent = upward)
\draw[->, thick] (A) -- (H);
\draw[->, thick] (B) -- (H);
\draw[->, thick] (C) -- (A);
\draw[->, thick] (D) -- (A);
\draw[->, thick] (E) -- (B);

% Labels on arrows
\node[left] at (-1.5, -1) {$\pointer{A_1}{H_0}$};
\node[right] at (1.5, -1) {$\pointer{B_1}{H_0}$};
\node[left] at (-2.8, -3) {$\pointer{C_2}{A_1}$};
\node[right] at (-0.3, -3) {$\pointer{D_2}{A_1}$};
\node[right] at (2.8, -3) {$\pointer{E_2}{B_1}$};

% Root annotation
\node[below=0.3cm of H] {\begin{tabular}{c}Root human\\$\pointer{H_0}{\bot}$\end{tabular}};

\end{tikzpicture}
\caption{Immutable spawn chain structure. Each directed edge $(c, p)$ represents $\pointer{c}{p}$. Arrows point from child to parent, upward toward the root. Pointers are sealed at spawn and never modified.}
\label{fig:spawn-chain}
\end{figure}

Figure~\ref{fig:spawn-chain} illustrates the chain structure for a representative spawn tree. Each entity carries a single immutable parent pointer (the sole exception being the root human, whose pointer is $\bot$). Traversing from any entity upward through its parent pointers reaches the root human in a finite number of steps. No back edges, cross edges, or cycles are possible due to the immutability constraint.

\subsection{Summary}

The immutable parent pointer is the structural foundation of the Recursive Spawning framework. Once assigned at spawn, $\texttt{ParentPointer}(c) = p$ becomes an indelible part of entity $c$'s record, enforced at the Fact Layer by the substrate itself. The $\texttt{Chain}(a)$ operator recursively constructs an entity's ancestry by following these pointers to the root, where $\texttt{ParentPointer}(root) = \bot$ terminates the ascent. Spawn authorization occurs exclusively at the Purpose Layer prior to Fact Layer instantiation, ensuring that every parent-child relationship reflects an authorized, recorded spawn event. The immutability guarantee propagates into the chain's non-repudiation, cycle-freedom, auditability, and temporal ordering — properties that the rest of the Recursive Spawning framework depends upon for authorization, inheritance, liability, and governance operations.